% =====================================================================
% AeroNet Lite — Final Report
% BS Data Science, AI Semester Project (SP2026)
% Compiles cleanly on Overleaf (pdflatex). Place all PNG figures
% from report/figures/ into a sub-folder called `figures/` next to this
% file before compiling.
% =====================================================================
\documentclass[11pt,a4paper]{article}

\usepackage[T1]{fontenc}
\usepackage[utf8]{inputenc}
\usepackage[margin=1in]{geometry}
\usepackage{graphicx}
\usepackage{booktabs}
\usepackage{caption}
\usepackage{subcaption}
\usepackage{xcolor}
\usepackage{listings}
\usepackage{amsmath,amssymb}
\usepackage{enumitem}
\usepackage{float}
\usepackage{tcolorbox}
\usepackage{algorithm}
\usepackage[noend]{algpseudocode}
\usepackage{hyperref}
\hypersetup{
  colorlinks=true,
  linkcolor=black,
  urlcolor=blue!60!black,
  citecolor=black
}

% ----- Code listing style ---------------------------------------------
\definecolor{codebg}{HTML}{F7F7F7}
\definecolor{codekw}{HTML}{1F4E79}
\definecolor{codestr}{HTML}{8B0000}
\definecolor{codecmt}{HTML}{6A737D}
\lstdefinestyle{py}{
  language=Python,
  basicstyle=\ttfamily\footnotesize,
  keywordstyle=\color{codekw}\bfseries,
  stringstyle=\color{codestr},
  commentstyle=\color{codecmt}\itshape,
  backgroundcolor=\color{codebg},
  showstringspaces=false,
  breaklines=true,
  frame=single,
  rulecolor=\color{black!20},
  numbers=left,
  numberstyle=\tiny\color{black!50},
  numbersep=8pt,
  xleftmargin=2.5em,
  framexleftmargin=2em
}

% Coloured side note box
\newtcolorbox{notebox}[1][]{
  colback=blue!4,
  colframe=blue!50!black,
  boxrule=0.6pt,
  arc=2pt,
  left=6pt, right=6pt, top=4pt, bottom=4pt,
  fonttitle=\bfseries\small,
  #1
}

\graphicspath{{figures/}}

% =====================================================================
\title{
  \vspace{-1.5cm}
  \textbf{AeroNet Lite}\\[0.4em]
  \large Autonomous Drone Delivery Simulation with CSP, Fleet Planning,\\
  A* Routing, Real-Time Replanning, Demand Forecasting, and Anomaly Detection
}
\author{
  \textbf{Ifham Razi} \\ 23i-2550
  \and
  \textbf{Romaan Khawar} \\ 23i-2554
  \and
  \textbf{Ali Haroon} \\ 23i-2592
}
\date{BS Data Science --- AI Semester Project (SP2026) \\ \today}

\begin{document}
\maketitle

% =====================================================================
\begin{abstract}
\noindent
AeroNet Lite is a lightweight autonomous drone-delivery simulator that
integrates five core AI techniques on a shared $10\!\times\!10$ city grid:
Constraint Satisfaction (layout validation), heuristic and Genetic
Algorithm search (fleet selection), A* pathfinding (delivery routing),
online optimization (real-time replanning around no-fly cells), and
supervised learning (demand regression and flight-anomaly classification).
Three real datasets back the regressors --- Bike Sharing Demand
(10{,}886 rows), Amazon Delivery (43{,}551 rows), and US City Population
Densities --- while a deliberately-constructed synthetic telemetry set
backs the four-class anomaly classifier. The system runs as a 20-step
console simulation \emph{and} as an interactive Plotly-based Streamlit
dashboard with step-through, restart, and seed-reproducible operation.
A procedural grid generator produces a fresh, CSP-valid layout for every
run. All ML training is reproducible (\texttt{random\_state=42}), giving
identical confusion matrices and metrics across reruns, while scenario
seeds vary deliveries, anomalies, and disruptions per run.
\end{abstract}

\vspace{0.4em}
\noindent\hrulefill\par
\tableofcontents
\noindent\hrulefill

% =====================================================================
\section{Introduction}
\label{sec:intro}

\subsection{Motivation}
Modern logistics increasingly evaluates autonomous drone fleets for
last-mile delivery, particularly for time-critical medical supplies. A
full city-scale drone simulator with realistic physics, traffic, weather,
and air-traffic-control integration is far outside the scope of a
two-week semester project. The original course brief explicitly maps
each ambitious sub-problem to a smaller analogue: from a full city to a
$10\!\times\!10$ grid, from ACO/GA over heavy modules to a brute-force
fleet search and four CSP rules, from a real-time UI to a static
matplotlib plus interactive Streamlit dashboard. The chosen scope still
exercises every AI technique required by the rubric --- CSP, search,
optimization, regression, and classification --- and integrates them
through a single shared grid-state contract.

\subsection{Project Goals}
The deliverables targeted in this project are:
\begin{itemize}[leftmargin=*]
  \item A modular Python project where every module is independently
        testable and integrated through a shared 10$\times$10 grid model.
  \item A 20-step console simulation matching the spec's example event
        log structure.
  \item An interactive Streamlit dashboard with an animated Plotly city
        map, drone trails, and live metrics.
  \item Five quantitative outputs supporting the rubric: a layout-validity
        report, a fleet selection under budget, A* routes, regression
        error metrics, and a classifier confusion matrix.
  \item Reproducibility separation: model training is pinned, but
        scenario seeds are configurable so demos can be varied without
        breaking model evaluation.
\end{itemize}

\subsection{Report Layout}
Section~\ref{sec:arch} describes the system architecture and module
contracts. Sections~\ref{sec:grid}--\ref{sec:ml} cover the five AI
modules in order. Sections~\ref{sec:data} and~\ref{sec:scenario} cover
the datasets used and the 20-step demonstration scenario.
Section~\ref{sec:dash} describes the Streamlit dashboard.
Sections~\ref{sec:challenges}--\ref{sec:limits} discuss implementation
challenges and limitations. Section~\ref{sec:contrib} attributes work to
each member.

% =====================================================================
\section{System Architecture}
\label{sec:arch}

\subsection{Module Contract}
Every module communicates through a single shared grid model, which
prevents the sub-systems from drifting out of sync as the simulation
progresses. Table~\ref{tab:modules} lists the read/write boundaries.

\begin{table}[H]
\centering
\caption{Modules, AI technique, and shared-state interaction.}
\label{tab:modules}
\begin{tabular}{@{}llp{4cm}p{4.5cm}@{}}
\toprule
\textbf{Module} & \textbf{AI Technique} & \textbf{Reads} & \textbf{Writes} \\
\midrule
Layout Validator     & CSP                          & zone, hub, charging flags        & validation results \\
Fleet Selector       & Heuristic / GA               & hubs, demand estimates           & selected drone list \\
Path Planner         & A* search                    & no-fly, costs, hubs, deliveries  & routes \\
Disruption Handler   & Online optimization          & active routes, no-fly flags      & rerouted segments \\
ML Pipeline          & Regression + Classification  & demand history, telemetry        & forecasts, anomaly labels \\
\bottomrule
\end{tabular}
\end{table}

The data flow is summarised in Figure~\ref{fig:arch}.

\begin{figure}[H]
\centering
\includegraphics[width=0.96\linewidth]{architecture_diagram.png}
\caption{High-level system architecture. AI modules (coloured blocks) read
from and write to the shared grid model in the centre. External datasets
(white boxes) feed the ML pipeline once at startup; their predictions
flow back to the grid as enriched densities, ETAs, and anomaly labels.}
\label{fig:arch}
\end{figure}

\subsection{Project Layout}
The repository is organised so each module has exactly one file, making
cross-referencing in viva straightforward.

\begin{lstlisting}[basicstyle=\ttfamily\footnotesize, frame=single, backgroundcolor=\color{codebg}, language=]
aeronet_lite/
  data/raw/                 (5 datasets: bike_sharing, amazon_delivery, ...)
  src/
    grid_model.py           shared 10x10 city + procedural generator
    layout_validator.py     CSP rules R1-R4
    fleet_selector.py       brute force AND genetic algorithm
    astar_planner.py        A* with commercial-corridor discount
    delivery_simulator.py   drone state machine, no-fly handling
    ml_pipeline.py          demand, anomaly, delivery-time models
    visualization.py        matplotlib report figures
    simulation.py           stepwise orchestrator
    main.py                 CLI entry point (--seed N for reproducibility)
    streamlit_app.py        interactive dashboard
  notebooks/                ML experimentation notebooks
  report/                   LaTeX source + figures + figure generator
\end{lstlisting}

% =====================================================================
\section{Shared Grid Model}
\label{sec:grid}

\subsection{Cell Schema}
Each cell in the $10\!\times\!10$ grid is represented as a Python
dataclass with the following fields:
\begin{lstlisting}[style=py]
@dataclass
class Cell:
    row: int
    col: int
    zone: Zone               # Residential/Commercial/Hospital/School/Industrial/Open
    density: int             # population density (sampled from real US data)
    is_hub: bool             # drone hub
    is_charging: bool        # charging pad
    is_medical_pickup: bool  # medical-supply pickup point
    no_fly: bool             # cell blocked for flight
    demand: float            # estimated delivery demand
\end{lstlisting}

\subsection{Procedural Grid Generation}
Rather than hard-coding a single example city, AeroNet Lite ships with a
constructive grid generator (\texttt{build\_random\_grid(seed)}) that
produces a fresh layout for every run while honouring all four CSP rules
\emph{by construction}. Algorithm~\ref{alg:gen} sketches the procedure.

\begin{algorithm}[H]
\caption{\textsc{BuildRandomGrid}(seed)}
\label{alg:gen}
\begin{algorithmic}[1]
\State $rng \gets \text{random}(seed)$
\State Place 2--3 hubs at mutual Manhattan distance $\geq 4$
\For{each hub $h$}
  \State Attach a charging pad at Manhattan distance $1$ \Comment{R3 holds}
\EndFor
\For{each hub $h$}
  \State Mark 4--7 nearby cells (within distance $3$) as Residential \Comment{R2 holds}
\EndFor
\State Grow a 2--3 cell Hospital cluster
\State Grow a 2--3 cell School cluster
\State Place a 2--4 cell Industrial cluster in a corner with no Hospital/School neighbour \Comment{R1 holds}
\State Scatter 8--14 Commercial cells through remaining open space
\State Mark a Medical Pickup adjacent to a Hospital \Comment{R4 holds}
\State Validate the layout; on failure retry with a perturbed sub-seed
\end{algorithmic}
\end{algorithm}

The four CSP rules from Section~\ref{sec:csp} hold by construction
because each placement rule is intersected with the relevant constraint
before commit. A 20-seed stress test confirmed 0/20 fall-backs to the
hardcoded sample grid. Figure~\ref{fig:gallery} shows four representative
layouts.

\begin{figure}[H]
\centering
\includegraphics[width=0.92\linewidth]{grid_gallery.png}
\caption{Four procedurally generated valid layouts. Hubs (\texttt{H},
blue squares), charging pads (\texttt{C}, teal squares), medical pickup
(red \texttt{+}), and any pre-placed no-fly cells appear as overlays.
Each grid satisfies R1--R4 by construction.}
\label{fig:gallery}
\end{figure}

\subsection{Population-Density Enrichment}
The cells' \texttt{density} field is then sampled from real
US-city-density data (Section~\ref{sec:data}). The 33-percentile
buckets are mapped to zone types so that residential cells draw from the
upper tertile, hospitals/schools from the middle, and industrial/open
from the lower tertile. The sampling is also seed-driven, so the same
\texttt{--seed} reproduces both the layout \emph{and} the density
distribution.

\begin{figure}[H]
\centering
\includegraphics[width=0.62\linewidth]{zone_map.png}
\caption{Single sample run rendered through the matplotlib visualization
helper. The visual identity (zone palette + landmark glyphs) is preserved
in the Streamlit dashboard.}
\label{fig:zone}
\end{figure}

% =====================================================================
\section{Module 1 --- Layout Validation (CSP)}
\label{sec:csp}

\subsection{Constraint Set}
The validator behaves as a CSP constraint checker over four rules. All
violations are accumulated rather than failing on the first error so the
report can suggest fixes for every issue found.

\begin{table}[H]
\centering
\caption{CSP rules and the geometric check used.}
\begin{tabular}{@{}llp{8.5cm}@{}}
\toprule
\textbf{ID} & \textbf{Rule} & \textbf{Check} \\
\midrule
R1 & Industrial $\not\!\!\sim$ School/Hospital   & For each Industrial cell, no 4-neighbour is School or Hospital. \\
R2 & Residential within 3 of a hub              & $\min_{h\in\text{hubs}} \|c-h\|_1 \leq 3$ for every Residential cell. \\
R3 & Hub within 2 of a charging pad             & $\min_{p\in\text{pads}} \|h-p\|_1 \leq 2$ for every Hub. \\
R4 & At least one Hospital within 1 of medical pickup & $\exists h \in \text{Hospitals}, m \in \text{Medical}: \|h-m\|_1 \leq 1$. \\
\bottomrule
\end{tabular}
\end{table}

\subsection{Implementation}
Each rule is implemented as a pure function returning a list of
\texttt{Violation} objects. The aggregator collects them all so the
report can list every offending cell with a suggested fix.

\begin{lstlisting}[style=py]
def check_residential_coverage(grid):
    violations = []
    hubs_list = hubs(grid)
    for cell in cells_of(grid, lambda c: c.zone == Zone.RESIDENTIAL):
        if not any(manhattan(cell.pos, h.pos) <= 3 for h in hubs_list):
            violations.append(Violation(
                rule_id="R2",
                cell=cell.pos,
                message="Residential cell more than 3 cells from any hub",
                suggestion="Add a hub closer to this district or convert "
                           "the cell to Open Field.",
            ))
    return violations
\end{lstlisting}

\begin{notebox}
\textbf{Why a CSP framing?} The rules could be encoded as a single
post-hoc assertion, but treating each rule as an independent constraint
function lets the procedural generator and validator share the same
declarative semantics --- a constraint that the generator enforces is
also one the validator checks, with no risk of drift.
\end{notebox}

% =====================================================================
\section{Module 2 --- Fleet Selection}
\label{sec:fleet}

\subsection{Drone Catalogue and Fitness}
The fleet selector chooses a small drone fleet under a fixed budget of
\$10{,}000 from two drone types:

\begin{table}[H]
\centering
\caption{Drone catalogue.}
\begin{tabular}{@{}lcccp{4cm}@{}}
\toprule
\textbf{Type} & \textbf{Cost (\$)} & \textbf{Payload} & \textbf{Range (cells)} & \textbf{Use case} \\
\midrule
Light & 1{,}000 & 2 kg & 12 & Short residential/commercial deliveries \\
Heavy & 1{,}800 & 5 kg & 20 & Long-range medical packages \\
\bottomrule
\end{tabular}
\end{table}

The fitness function rewards demand coverage and penalises overspending:
\begin{equation}
\text{score}(\ell, h) =
\begin{cases}
0.75 \cdot \text{coverage}_{\%}(\ell, h) - 0.25 \cdot \text{budget\_used}_{\%}(\ell, h) & \text{if cost} \leq B \\
-\infty & \text{otherwise}
\end{cases}
\label{eq:fitness}
\end{equation}
where $\ell$ is the light count, $h$ the heavy count, and $B = 10{,}000$.

\subsection{Two Algorithms}
We implemented both algorithms permitted by the spec for AI breadth:

\paragraph{Brute force.} Enumerate every $(\ell, h) \in [0, 10]^2$ and
keep the highest-scoring feasible plan. With only 121 candidates this is
microsecond-fast and provides a ground-truth optimum for benchmarking
the GA.

\paragraph{Genetic Algorithm.} A small evolutionary search demonstrates
the rubric's "evolutionary search" requirement.

\begin{algorithm}[H]
\caption{\textsc{FleetGA}(grid, budget, generations $=30$, population $=20$)}
\label{alg:ga}
\begin{algorithmic}[1]
\State Seed RNG with $42$
\State $P \gets$ population of 20 random $(\ell, h)$ chromosomes
\State $best \gets \arg\max_{\chi \in P} \text{fitness}(\chi)$
\For{$g \gets 1$ to generations}
  \State $children \gets \emptyset$
  \While{$|children| < |P|$}
    \State $p_1, p_2 \gets$ tournament-select($P$, $k=3$) twice
    \State $\chi \gets$ single-point crossover$(p_1, p_2)$
    \State with prob.\ $0.20$ mutate one gene by $\pm 1$
    \State append $\chi$ to $children$
  \EndWhile
  \State $P \gets children \cup \{best\}$ \Comment{elitism}
  \State update $best$ if any $\chi \in P$ scores higher
\EndFor
\State \Return $best$
\end{algorithmic}
\end{algorithm}

\subsection{Result}
On the typical generated grid both algorithms converge to
\textbf{5 Light + 2 Heavy drones}: cost \$8{,}600 (86\,\% of budget),
coverage 100\,\%, fitness $= 0.75 \times 100 - 0.25 \times 86 = 53.5$.
The GA matches the brute-force optimum within 5 generations.

% =====================================================================
\section{Module 3 --- A* Path Planner}
\label{sec:astar}

\subsection{Design Choices}
\begin{table}[H]
\centering
\begin{tabular}{@{}ll@{}}
\toprule
\textbf{Item} & \textbf{Choice} \\
\midrule
State          & grid coordinate $(r,c)$ \\
Actions        & up, down, left, right \\
Path cost      & $1.0$ default, $0.8$ when entering a Commercial cell \\
Blocked        & cells with \texttt{no\_fly == True} \\
Heuristic      & Manhattan distance $\|s - g\|_1$ (admissible) \\
Tiebreaker     & monotonic counter for stable ordering in \texttt{heapq} \\
\bottomrule
\end{tabular}
\end{table}

\subsection{Pseudocode}
\begin{algorithm}[H]
\caption{\textsc{A*}(start, goal, grid)}
\label{alg:astar}
\begin{algorithmic}[1]
\State $g \gets \{start: 0\}$, $parent \gets \{start: \text{None}\}$, $closed \gets \emptyset$
\State $open \gets$ heap with one entry $(h(start, goal), 0, start)$
\While{open is non-empty}
  \State $(f, \_, n) \gets$ pop min from $open$
  \If{$n = goal$} \Return reconstruct($parent, n$) \EndIf
  \State add $n$ to $closed$
  \For{each $n' \in$ neighbours($n$) with \texttt{not no\_fly}}
    \If{$n' \in closed$} \textbf{continue} \EndIf
    \State $g' \gets g[n] + \texttt{move\_cost}(grid, n')$
    \If{$g' < g.get(n', \infty)$}
      \State $g[n'] \gets g'$, $parent[n'] \gets n$
      \State push $(g' + h(n', goal),\, \text{counter}, n')$ onto $open$
    \EndIf
  \EndFor
\EndWhile
\State \Return failure
\end{algorithmic}
\end{algorithm}

\subsection{Three-Segment Delivery Routing}
A delivery is decomposed into three A* segments:
$\text{hub} \to \text{pickup} \to \text{dropoff} \to \text{hub}$. Failed
segments mark the delivery as \emph{unroutable} so the simulator can log
the failure rather than crash. Figure~\ref{fig:routes} shows the route
state at the start of a run and after the disruption handler has fired.

\begin{figure}[H]
\centering
\begin{subfigure}[t]{0.48\linewidth}
  \centering
  \includegraphics[width=\linewidth]{routes_initial.png}
  \caption{Initial planned routes (step 6).}
\end{subfigure}
\hfill
\begin{subfigure}[t]{0.48\linewidth}
  \centering
  \includegraphics[width=\linewidth]{routes_final.png}
  \caption{Routes after rerouting around an activated no-fly cell.}
\end{subfigure}
\caption{Initial vs.\ post-disruption A* routes for the active fleet.}
\label{fig:routes}
\end{figure}

\subsection{Cost Model Rationale}
The 0.8 commercial-corridor discount is intended to model preferred
flight corridors over commercial districts (where there is less risk of
flying over people). Without it A* paths just go diagonally; with it,
drones realistically prefer the wider commercial axes between hubs.
Manhattan distance remains an admissible heuristic in the discounted case
because every movement still costs at least 0.8, so $h \leq g^*$ holds.

% =====================================================================
\section{Module 4 --- Real-Time Disruption Handling}
\label{sec:disrupt}

\subsection{No-Fly Activation}
At step 11 of each run, the simulator activates a no-fly cell on a cell
that lies a few steps ahead of an active drone (so the rerouting
demonstration is non-trivial rather than firing in empty space). The
handler iterates over every active drone and, when the no-fly intersects
\emph{any} of its remaining segments (current \emph{or} future), replans
all remaining segments from the drone's current position via repeated A*
calls. Algorithm~\ref{alg:reroute} sketches the procedure.

\begin{algorithm}[H]
\caption{\textsc{Reroute}(drone, grid)}
\label{alg:reroute}
\begin{algorithmic}[1]
\State $goals \gets$ list of segment endpoints from $drone$.seg\_idx onward
\State $start \gets drone$.current\_pos
\State $newSegments \gets$ already-completed segments
\For{each $goal \in goals$}
  \State $r \gets$ \textsc{A*}($start, goal, grid$)
  \If{$r$.success is False}
    \State mark $drone$ delayed; mark delivery delayed
    \State \textbf{return}
  \EndIf
  \State append $r$.path to $newSegments$; $start \gets goal$
\EndFor
\State $drone$.segments $\gets newSegments$
\State $drone$.pos\_in\_seg $\gets 0$
\end{algorithmic}
\end{algorithm}

\subsection{Anomaly-Driven Disruption}
A second event hook fires at step 18: synthetic telemetry is generated
for a non-idle drone and classified by the anomaly model
(Section~\ref{sec:ml}). On any non-Normal label the drone is force-returned
to its hub via a fresh A* call from its current position, aborting its
in-flight delivery.

% =====================================================================
\section{Module 5 --- ML Pipeline}
\label{sec:ml}

The pipeline trains three models. All training uses
\texttt{random\_state=42} so results are bit-for-bit reproducible.

\begin{table}[H]
\centering
\caption{Trained models, datasets, and metrics.}
\label{tab:models}
\begin{tabular}{@{}lllll@{}}
\toprule
\textbf{Model}       & \textbf{Algorithm}        & \textbf{Dataset}     & \textbf{Rows}  & \textbf{Metric} \\
\midrule
Demand               & Random Forest             & Bike Sharing Demand   & 10{,}886       & MAE 56.55, RMSE 86.65 \\
Demand (baseline)    & Linear Regression         & Bike Sharing Demand   & 10{,}886       & MAE 112.16, RMSE 152.78 \\
Delivery time        & RF + One-Hot Pipeline     & Amazon Delivery       & 43{,}551       & MAE 26.31\,min, RMSE 37.67\,min \\
Anomaly classifier   & Random Forest (4-class)   & Synthetic telemetry   & 2{,}000        & Accuracy 1.00 \\
\bottomrule
\end{tabular}
\end{table}

\subsection{Demand Forecasting}
\textbf{Features:} \texttt{hour, day\_of\_week, temperature, weather}.
We train both a Linear Regression baseline and a Random Forest on a
20\,\% held-out split. The Random Forest wins decisively on both error
metrics (Figure~\ref{fig:demand_compare}) and is selected as the deployed
model.

\begin{figure}[H]
\centering
\includegraphics[width=0.78\linewidth]{demand_model_comparison.png}
\caption{Linear Regression baseline vs.\ Random Forest on the Bike
Sharing Demand dataset (10{,}886 rows). Random Forest cuts MAE by
$\sim 50$\,\% and RMSE by $\sim 43$\,\%, justifying its selection over
the simpler model.}
\label{fig:demand_compare}
\end{figure}

The trained model drives the per-cell demand prediction shown in
Figure~\ref{fig:heatmap}.

\begin{figure}[H]
\centering
\includegraphics[width=0.62\linewidth]{demand_heatmap.png}
\caption{Predicted per-cell demand at hour 12, weekday, $20\,^\circ\!$C,
clear weather, scaled by the cell's enriched population density.}
\label{fig:heatmap}
\end{figure}

\subsection{Feature Importance}
Random Forest exposes per-feature Gini importance after training. For
the demand model, hour-of-day dominates --- consistent with the
twin morning/evening commute peaks in the bike-sharing data. For the
anomaly model, all four engineered features carry comparable weight,
which is consistent with the synthetic data design where each anomaly
class is driven by its own dominant feature (Figure~\ref{fig:fimp}).

\begin{figure}[H]
\centering
\includegraphics[width=0.96\linewidth]{feature_importance.png}
\caption{Random Forest feature importances. Left: demand model on Bike
Sharing data --- hour-of-day dominates. Right: anomaly classifier on
synthetic telemetry --- spread across the four engineered features that
distinguish each anomaly class.}
\label{fig:fimp}
\end{figure}

\subsection{Delivery-Time Prediction}
A scikit-learn \texttt{Pipeline} combining \texttt{ColumnTransformer}
+ \texttt{OneHotEncoder} + Random Forest predicts ETA in minutes from
numeric features (\texttt{distance\_km, agent\_age, agent\_rating}) and
categorical features (\texttt{weather, traffic, vehicle, area}). The
distance is computed via haversine from store and drop-off coordinates.
Figure~\ref{fig:eta} compares predictions against ground truth on the
held-out test split.

\begin{figure}[H]
\centering
\includegraphics[width=0.65\linewidth]{delivery_time_scatter.png}
\caption{Amazon Delivery time regressor: actual vs.\ predicted minutes
on a 20\,\% held-out split (n $=$ 8{,}710). Predictions cluster along
the $y=x$ ideal line; MAE $=$ 26.31\,min, RMSE $=$ 37.67\,min.}
\label{fig:eta}
\end{figure}

The simulator uses this model to annotate every assigned delivery with a
baseline ETA, scaled by the assumption that one Manhattan grid cell
$\approx 0.5$\,km of real-world distance.

\subsection{Anomaly Classification}
Synthetic telemetry produces four labelled regimes by construction:
\begin{itemize}[leftmargin=*]
  \item \textbf{Normal} --- low battery drop and route deviation.
  \item \textbf{BatteryAnomaly} --- unusually high battery drop.
  \item \textbf{RouteAnomaly} --- unusually high route deviation.
  \item \textbf{SensorSpike} --- high altitude or speed change.
\end{itemize}
With clean class boundaries the classifier reaches 100\,\% accuracy on a
stratified 20\,\% test split. The matrix is fully deterministic across
runs because the synthetic data, the train/test split, and the Random
Forest all share \texttt{random\_state=42}.

\begin{figure}[H]
\centering
\includegraphics[width=0.55\linewidth]{confusion_matrix.png}
\caption{Anomaly-classifier confusion matrix on the 400-row test split.
The clean diagonal (zero off-diagonal mass) is a property of the
deliberately separable synthetic data, not an artefact of overfitting,
and is acknowledged explicitly in the limitations section.}
\label{fig:cm}
\end{figure}

% =====================================================================
\section{Datasets}
\label{sec:data}

Three real datasets feed the pipeline; the project treats anomaly data
as synthetic per the course guidance, since real UAV telemetry (CMU ALFA
dataset) requires bespoke ROS-bag processing beyond the two-week budget.

\begin{table}[H]
\centering
\caption{Datasets used in the project.}
\begin{tabular}{@{}llp{6cm}@{}}
\toprule
\textbf{Purpose} & \textbf{Source} & \textbf{Role in the system} \\
\midrule
Demand forecasting   & Kaggle Bike Sharing Demand            & Hourly demand proxy for the regression model. \\
Delivery time        & Kaggle Amazon Delivery                & ETA regressor with realistic feature mix. \\
Grid density         & Kaggle US City Population Densities   & Bucketed Low/Medium/High densities sampled into cells per zone. \\
Anomaly (chosen)     & Synthetic                              & 4-class clean separation; documented assumption. \\
Anomaly (alternative)& CMU ALFA UAV Failure Dataset           & Available on disk; not wired (per-flight ROS topics). \\
\bottomrule
\end{tabular}
\end{table}

\subsection{Auto-Detection and Fallback}
The pipeline checks for each CSV at a fixed location under
\texttt{data/raw/}; if a file is missing the demand model gracefully
falls back to in-process synthetic data so the demo always runs. The
delivery-time and population-density loaders are silent no-ops on
missing data, which simplifies setup for graders without download
access.

% =====================================================================
\section{20-Step Simulation Scenario}
\label{sec:scenario}

The simulation is structured as 20 logical steps so the demonstration
covers all modules. Each step prints to an event log and updates the
visual state.

\begin{table}[H]
\centering
\caption{Step-by-step scenario plan.}
\begin{tabular}{@{}clp{8.5cm}@{}}
\toprule
\textbf{Step}  & \textbf{Phase}             & \textbf{Action} \\
\midrule
1--3   & Init / Validate / Fleet     & Generate grid, run CSP validator, brute-force fleet selection. \\
4--6   & ML training + Deliveries    & Train demand, anomaly, delivery-time models; assign 4 deliveries with ETAs. \\
7--10  & Movement                    & Drones advance two cells per step along their A* paths. \\
11     & Disruption                  & Activate a no-fly cell on an active route; replan affected drones. \\
12--14 & Movement (post-disruption)  & Continue along the new paths. \\
15--17 & Forecast + Extra delivery   & Run demand forecast; assign an additional delivery to an idle drone. \\
18     & Anomaly injection           & Synthesize anomalous telemetry; classify; flag the affected drone. \\
19     & Anomaly response            & Force return-to-hub for non-Normal classifications. \\
20     & Final tick + Summary        & Print completed / delayed / failed / in-progress counts. \\
\bottomrule
\end{tabular}
\end{table}

\begin{figure}[H]
\centering
\includegraphics[width=0.85\linewidth]{event_timeline.png}
\caption{Event log produced during a typical 20-step run.}
\label{fig:timeline}
\end{figure}

\paragraph{Reproducibility.} The CLI accepts an optional
\texttt{--seed N} argument. Each run prints the seed used so any output
shown in this report can be reproduced with
\texttt{python -m src.main --seed N}.

% =====================================================================
\section{Interactive Streamlit Dashboard}
\label{sec:dash}

The dashboard wraps the same \texttt{Simulation} class used by the CLI
but exposes per-step advancement, restart, and seeded replay through a
browser UI.

\subsection{Visualisation Stack}
The map is rendered with Plotly so cells, drones, routes, and tooltips
are all interactive:
\begin{itemize}[leftmargin=*]
  \item \textbf{City map} --- pastel-coloured cells per zone with
        building emojis ($\mathrm{\unicode{x1F3E0}}, \mathrm{\unicode{x1F3EC}},
        \mathrm{\unicode{x1F3E5}}, \mathrm{\unicode{x1F3EB}},
        \mathrm{\unicode{x1F3ED}}, \mathrm{\unicode{x1F33F}}$),
        HUB / CHG / MED overlays, drone triangles with status-coloured
        outlines, dotted planned routes, solid trails, and an X overlay
        for active no-fly cells. Drones at the same hub are arranged in
        a small ring so each is independently visible.
        Hovering on any cell or drone shows a rich tooltip.
  \item \textbf{Controls} --- Restart (re-run the same seed), New
        random (fresh seed), Step (advance one), Run to end, Auto-play
        with a speed slider.
  \item \textbf{Tabs} --- Event log, Drone status, Delivery table,
        Demand heatmap, Model metrics, Layout report.
\end{itemize}

\subsection{Caching}
The three ML models are wrapped in
\texttt{@st.cache\_resource}, so they are trained once per Streamlit
session and reused across every restart. This makes the
\textit{Restart} button effectively instant after the first warm-up.

\subsection{Launching}
The dashboard is launched with:
\begin{lstlisting}[language=bash, basicstyle=\ttfamily\footnotesize, frame=single, backgroundcolor=\color{codebg}]
cd d:\codes\ai2
.\.venv\Scripts\python.exe -m streamlit run src\streamlit_app.py
\end{lstlisting}

% =====================================================================
\section{Implementation Challenges}
\label{sec:challenges}

\subsection{State Machine Bugs in the Drone Lifecycle}
The first end-to-end run revealed an \texttt{IndexError} on the final
step: when \texttt{force\_return\_to\_hub} replaces a drone's three-leg
plan with a single segment, the original step\_drone routine still
expected exactly three segments. The fix was to use
\texttt{len(drone.segments)} instead of the literal \texttt{3} as the
end-of-plan condition.

A subtler bug followed: a drone that landed on the last cell of the
last segment was not marked complete until the next tick, so the
20-step demo summary showed deliveries that were geometrically finished
as still in progress. The fix was to detect arrival on the final cell of
the final segment within the \emph{same} call and trigger completion
immediately.

\subsection{No-Fly Activation Did Not Trigger Reroutes}
Initially the no-fly handler only inspected the current segment of each
drone. A drone whose pickup landed at the same cell as a downstream
no-fly point would proceed normally to the pickup, then route through
the no-fly cell on the next segment because that segment had been
planned before the disruption. The fix was to inspect \emph{all
remaining} segments and, when any of them crosses the new no-fly cell,
re-plan from the drone's current position through every remaining
segment goal.

\subsection{Demand Heatmap Aspect Ratio}
The Plotly heatmap initially used \texttt{scaleanchor="x"} without a
constrained x range, which caused it to fill the entire wide container
while the y axis stayed locked to the data span. The fix was to set
explicit ranges on both axes plus \texttt{constrain="domain"}, which
forces square cells regardless of container width.

\subsection{Drones Stacked at the Same Hub}
When multiple drones share a home hub (typical with brute-force fleet
selection), all their markers stacked in the same screen position so
only one was visible. The fix is a small ring-offset layout: drones at
the same cell are placed at evenly-spaced angles around a 0.28-cell
radius from the cell centre.

\subsection{Procedural Grid Generator's CSP Compliance}
The first generator placed cells uniformly at random and then validated
the result, which gave a high failure rate (R2 violations were the
common offender). Switching to a constructive approach where each
placement step intersects with the relevant constraint gave 0/20
failures across the seeded stress test.

% =====================================================================
\section{Results and Evaluation}
\label{sec:results}

\paragraph{Layout validator.} On the procedurally generated grid, all
four CSP rules pass; on the spec's deliberately-broken sample grid the
validator correctly reports R2 and R3 violations with offending cell
coordinates and human-readable suggested fixes.

\paragraph{Fleet.} Both brute force and the GA converge to the same
solution: 5 Light + 2 Heavy, cost \$8{,}600, coverage 100\,\%, fitness
$53.5$.

\paragraph{Routing and disruption.} A* finds a feasible route for every
delivery on the procedurally generated grids tested. Activating a no-fly
cell on an active drone's path triggers a successful reroute in every
test seed evaluated (20 distinct grids, 4 deliveries each).

\paragraph{Models.} Bike Sharing demand: RF MAE $= 56.55$, RMSE $= 86.65$
versus LR baseline MAE $= 112.16$, RMSE $= 152.78$; the RF is selected.
Amazon Delivery time: RF MAE $= 26.31$ min, RMSE $= 37.67$ min on a
43{,}551-row split. Anomaly classifier: 100\,\% accuracy on the synthetic
test split with the confusion matrix shown in Figure~\ref{fig:cm}.

\paragraph{End-to-end demo.} A typical 20-step run produces 2 completed,
1 delayed (anomaly-aborted), 0 failed, and 2 still-in-progress deliveries.
Different seeds produce different layouts, deliveries, and anomaly types
while the underlying models stay deterministic (Section~\ref{sec:ml}).

\paragraph{Rubric alignment.}
\begin{table}[H]
\centering
\begin{tabular}{@{}p{4cm}p{10cm}@{}}
\toprule
\textbf{Rubric area}      & \textbf{Where it is satisfied} \\
\midrule
Design document            & Sections~\ref{sec:arch}--\ref{sec:dash}: per-module algorithm choice, data model, UI plan. \\
Technical implementation   & Working Python project; each module standalone and integrated through the shared grid. \\
AI concept coverage        & CSP (\S\ref{sec:csp}), Search/GA (\S\ref{sec:fleet}), A* (\S\ref{sec:astar}), Optimization (\S\ref{sec:disrupt}), Regression \& Classification (\S\ref{sec:ml}). \\
Viva and defense           & All algorithms are first-principles implementations explainable from this report. \\
Interface and presentation & Static matplotlib figures (this report) + interactive Plotly + Streamlit dashboard (\S\ref{sec:dash}). \\
\bottomrule
\end{tabular}
\end{table}

% =====================================================================
\section{Limitations}
\label{sec:limits}

\begin{itemize}[leftmargin=*]
  \item \textbf{Anomaly classifier accuracy is "too clean."} The
        synthetic data was designed to cleanly separate the four classes,
        so the 100\,\% test accuracy reflects the data design rather than
        a generalising model. We document this explicitly to avoid
        misleading the viva.
  \item \textbf{No multi-drone collision avoidance.} Routes are planned
        independently per drone; two drones can in principle occupy the
        same cell at the same step. A real system would need temporal A*.
  \item \textbf{Static cost model.} A* uses a fixed 1.0 / 0.8 cost; we do
        not model wind, weight-dependent battery drain, or dynamic
        congestion.
  \item \textbf{Deterministic ML training.} For reproducibility we pin
        \texttt{random\_state=42}. We did not run cross-validation or
        bootstrapping to report confidence intervals on the metrics.
  \item \textbf{Real UAV telemetry not wired.} The CMU ALFA dataset is
        present on disk but not loaded into the anomaly classifier --- it
        is per-flight ROS topics that need bespoke processing.
  \item \textbf{Single-city scope.} The grid is fixed at $10\!\times\!10$;
        nothing currently scales to a larger environment without
        re-tuning hub spacing and budget.
\end{itemize}

% =====================================================================
\section{Per-Member Contributions}
\label{sec:contrib}

The work was split across the team along the module boundaries described
in Section~\ref{sec:arch}. All members participated in design discussions
and the final report.

\begin{table}[H]
\centering
\caption{Contribution split.}
\begin{tabular}{@{}p{4cm}p{10cm}@{}}
\toprule
\textbf{Member}        & \textbf{Primary contributions} \\
\midrule
Ifham Razi (23i-2550)  & Shared grid model, procedural grid generator,
                         CSP validator, integration through
                         \texttt{Simulation} class, dashboard layout. \\
\addlinespace
Romaan Khawar (23i-2554) & A* path planner, three-segment delivery
                          routing, real-time disruption handler,
                          Plotly city renderer, drone state machine. \\
\addlinespace
Ali Haroon (23i-2592)   & Fleet selector (brute force + GA), ML pipeline
                         (demand, anomaly, delivery-time models), feature
                         engineering, dataset wiring, report figures. \\
\bottomrule
\end{tabular}
\end{table}

% =====================================================================
\section{Future Work}
\label{sec:future}

\begin{itemize}[leftmargin=*]
  \item Wire the CMU ALFA telemetry into the anomaly classifier to
        replace the synthetic dataset and report a more realistic
        confusion matrix.
  \item Replace the Manhattan A* heuristic with a wind/airspace-aware
        cost once altitude is added to the cell model.
  \item Add multi-drone collision avoidance by extending A* to a
        space-time graph with reservation tables.
  \item Smooth-frame Plotly animation between simulation steps for the
        dashboard, using \texttt{frames} so drone movement is
        interpolated rather than discrete.
  \item Train the demand model with a wider feature set
        (humidity, wind speed, holiday flags) and report cross-validated
        confidence intervals.
\end{itemize}

% =====================================================================
\section{Conclusion}
\label{sec:conclusion}

AeroNet Lite delivers all five required AI components in an integrated
$10\!\times\!10$ simulator with both a deterministic CLI and an
interactive dashboard. Real datasets back the regression models;
synthetic data backs the anomaly classifier with a clearly-documented
assumption. The procedural grid generator and seeded runs make the system
both reproducible and varied --- two properties that often pull against
each other --- by separating model-training seeds (pinned) from
scenario seeds (configurable). The codebase totals roughly 2{,}000 lines
of Python, organised so each module can be opened, read, and explained
from first principles in viva.

% =====================================================================
\section*{References}
\label{sec:refs}
\addcontentsline{toc}{section}{References}

\begin{enumerate}[leftmargin=*]
  \item Kaggle --- Bike Sharing Demand. \url{https://www.kaggle.com/c/bike-sharing-demand}
  \item Kaggle --- Amazon Delivery Dataset. \url{https://www.kaggle.com/datasets/sujalsuthar/amazon-delivery-dataset}
  \item Kaggle --- US City Population Densities. \url{https://www.kaggle.com/datasets/mmcgurr/us-city-population-densities}
  \item CMU KiltHub --- ALFA UAV Fault and Anomaly Detection Dataset. \url{https://kilthub.cmu.edu/articles/dataset/ALFA_A_Dataset_for_UAV_Fault_and_Anomaly_Detection/12707963}
  \item Hart, P., Nilsson, N., Raphael, B. (1968). \emph{A Formal Basis for the Heuristic Determination of Minimum Cost Paths}. IEEE Transactions on Systems Science and Cybernetics.
  \item Russell, S., Norvig, P. \emph{Artificial Intelligence: A Modern Approach}, 4th ed. (Constraint Satisfaction, Search, Machine Learning chapters.)
  \item Pedregosa, F., et al. (2011). \emph{Scikit-learn: Machine Learning in Python}. Journal of Machine Learning Research, 12, 2825--2830.
  \item Plotly Technologies Inc. (2015). \emph{Collaborative data science}. Plotly. \url{https://plot.ly}
  \item Streamlit Inc. \emph{Streamlit Documentation}. \url{https://docs.streamlit.io}
\end{enumerate}

\end{document}
